% !TEX root = ../manuscript.tex

\section{Field-Scale Multiphase Simulation and Analysis}

\subsection{Reservoir Model and Injection Scenario}

High-performance multiphase simulations model megatonne-scale \coTwo{}
injection and subsequent migration over 1,000 years in the approximately
two-million-cell offshore Texas model. The model resolves the storage
reservoir, top seal, and three-dimensional growth fault so that migration can
be tracked within the reservoir and both along and across the fault.

\todo{Report the simulator and nonlinear-solver configuration, fluid model,
initial and boundary conditions, injection rate and duration, well controls,
rock and fluid properties outside the fault, grid statistics, timestep
strategy, convergence tolerances, and computing platform.}

\subsection{Fault-Property Mapping and Ensemble Design}

Each of the 162 geologic scenarios contributes 10 sampled three-dimensional
fault-property cases, giving
\begin{equation}
162\times10 = 1{,}620
\end{equation}
field-scale simulations. Every case contains spatially variable fault
permeability and porosity over the $6\times87$ fault representation, together
with saturation-region assignments for the selected capillary-pressure and
relative-permeability curves. The permeability case design is deterministic
for a given geology and case type, while draws within the prescribed sampling
pools vary reproducibly along strike.

\todo{Define the final reservoir-ready file schema, property-to-grid mapping,
saturation-region treatment, random-seed policy, restart protocol, and
simulation acceptance criteria.}

\subsection{Migration and Containment Quantities of Interest}

The simulation ensemble is evaluated using quantities of interest that link
fault behavior to storage security:
\begin{enumerate}
  \item the fraction of injected \coTwo{} retained within the intended storage
  reservoir;
  \item the fraction entering the fault;
  \item the fraction entering the top seal;
  \item the maximum upward migration distance; and
  \item the partitioning of \coTwo{} between dissolved and free phases.
\end{enumerate}

\todo{Provide exact spatial masks, denominators, evaluation times, phase and
component definitions, and numerical thresholds for every quantity of
interest. Distinguish instantaneous maxima from values evaluated at the end of
injection and after 1,000 years.}

\subsection{Ensemble Analysis}

The ensemble analysis compares outcomes across geologic scenarios, within-
window stochastic states, and cross-window coupling assumptions. This
hierarchical labeling enables both total predictive ranges and conditional
comparisons that identify which uncertainty source controls each migration or
containment metric.

\todo{Specify the statistical estimands, uncertainty-decomposition method,
matched comparisons among the 10 case types, treatment of failed simulations,
and confidence or sensitivity diagnostics.}
